
%% ====================== ปก ======================
\begin{frame}[plain]
  \maketitle
\end{frame}

%% ====================== ทบทวน ======================
\begin{frame}{ทบทวนสัปดาห์ที่ 1}
  \begin{itemize}
    \item \textbf{ประพจน์} = ข้อความที่ระบุค่าความจริง (T/F) ได้แน่นอน
    \item \textbf{ตัวเชื่อม 5 ตัว}: $\neg,\ \land,\ \lor,\ \rightarrow,\ \leftrightarrow$
  \end{itemize}
  \begin{center}\small
  \begin{tabular}{|c|c|c|c|c|c|c|}
    \hline
    \textbf{$p$} & \textbf{$q$} & \textbf{$\neg p$} & \textbf{$p \land q$} & \textbf{$p \lor q$} & \textbf{$p \rightarrow q$} & \textbf{$p \leftrightarrow q$} \\ \hline
    \TT & \TT & \FF & \TT & \TT & \TT & \TT \\ \hline
    \TT & \FF & \FF & \FF & \TT & \FF & \FF \\ \hline
    \FF & \TT & \TT & \FF & \TT & \TT & \FF \\ \hline
    \FF & \FF & \TT & \FF & \FF & \TT & \TT \\ \hline
  \end{tabular}
  \end{center}
  \footnotesize สัปดาห์นี้: นำตารางค่าความจริงไป \textbf{จำแนกประเภทประพจน์ และตรวจสอบความสมมูล}
\end{frame}

%% ====================== วัตถุประสงค์ ======================
\begin{frame}{วัตถุประสงค์การเรียนรู้ (สัปดาห์ที่ 2)}
  เมื่อเรียนจบสัปดาห์นี้ นักศึกษาจะสามารถ
  \begin{enumerate}
    \item จำแนก \textbf{สัจนิรันดร์ ข้อขัดแย้ง และประพจน์กลาง} ได้
    \item ตรวจสอบ \textbf{สมมูลทางตรรกศาสตร์} และใช้ \textbf{กฎพีชคณิตของประพจน์} ได้
  \end{enumerate}
\end{frame}

%% ====================== หัวข้อ ======================
\begin{frame}{หัวข้อที่จะศึกษาในสัปดาห์นี้}
  \tableofcontents
\end{frame}


%% =====================================================================
\section{สัจนิรันดร์ ข้อขัดแย้ง และประพจน์กลาง}
%% =====================================================================

\begin{frame}{จำแนกประพจน์เชิงประกอบ 3 ประเภท}
  \begin{block}{}
  \begin{tabular}{@{}p{0.32\textwidth} p{0.62\textwidth}@{}}
    \textbf{\textcolor{Tgreen}{สัจนิรันดร์}} (Tautology) & จริงทุกกรณี \\[0.4em]
    \textbf{\textcolor{Fred}{ข้อขัดแย้ง}} (Contradiction) & เท็จทุกกรณี \\[0.4em]
    \textbf{\textcolor{navy}{ประพจน์กลาง}} (Contingency) & จริงบ้างเท็จบ้าง ขึ้นกับค่าตัวแปร \\
  \end{tabular}
  \end{block}
  \vspace{0.5em}
  \begin{itemize}
    \item ความสัมพันธ์: นิเสธของสัจนิรันดร์ = ข้อขัดแย้ง และในทางกลับกัน
    \item ประพจน์ส่วนใหญ่ในชีวิตจริงเป็น \textbf{ประพจน์กลาง}
  \end{itemize}
\end{frame}

\begin{frame}{สัจนิรันดร์ (Tautology)}
  \begin{block}{นิยาม}
    ประพจน์เชิงประกอบที่มีค่าความจริงเป็น \textbf{จริงในทุกกรณี}
  \end{block}
  \vspace{0.3em}
  ตัวอย่างคลาสสิก $p \lor \neg p$ (กฎกีดกันกลาง, Law of Excluded Middle)
  \begin{columns}[T]
    \begin{column}{0.45\textwidth}
      \begin{center}
      \begin{tabular}{|c|c|c|}
        \hline \textbf{$p$} & \textbf{$\neg p$} & \textbf{$p\lor\neg p$} \\ \hline
        \TT & \FF & \TT \\ \hline
        \FF & \TT & \TT \\ \hline
      \end{tabular}
      \end{center}
    \end{column}
    \begin{column}{0.5\textwidth}
      ตัวอย่างสัจนิรันดร์อื่น ๆ
      \begin{itemize}
        \item $p \rightarrow p$
        \item $p \leftrightarrow p$
        \item $(p \land q) \rightarrow p$
      \end{itemize}
    \end{column}
  \end{columns}
  \vspace{0.4em}
  คอลัมน์ผลลัพธ์เป็น \TT\ ทุกแถว $\Rightarrow$ เป็นสัจนิรันดร์
\end{frame}

\begin{frame}{$(p\land q)\rightarrow p$ เป็นสัจนิรันดร์}
  \begin{center}
  \begin{tabular}{|c|c|c|c|}
    \hline
    \textbf{$p$} & \textbf{$q$} & \textbf{$p\land q$} & \textbf{$(p\land q)\rightarrow p$} \\ \hline
    \TT & \TT & \TT & \TT \\ \hline
    \TT & \FF & \FF & \TT \\ \hline
    \FF & \TT & \FF & \TT \\ \hline
    \FF & \FF & \FF & \TT \\ \hline
  \end{tabular}
  \end{center}
  \vspace{0.8em}
  คอลัมน์สุดท้ายเป็น \TT\ ทุกแถว จึงเป็นสัจนิรันดร์
  \par\medskip
  \begin{exampleblock}{การประยุกต์}
    หากต้องการพิสูจน์ว่า $P$ จริง อาจแสดงว่า $\neg P$ นำไปสู่ข้อขัดแย้ง
    $\Rightarrow$ $\neg P$ เท็จ $\Rightarrow$ $P$ จริง (\textbf{Proof by Contradiction})
  \end{exampleblock}
\end{frame}

\begin{frame}{ตัวอย่าง: พิสูจน์ว่าไม่มีจำนวนเต็มที่ใหญ่ที่สุด}
  \textbf{วิธีพิสูจน์โดยข้อขัดแย้ง}
  \begin{enumerate}
    \item สมมติในทางตรงกันข้ามให้ \textbf{มีจำนวนเต็มที่ใหญ่ที่สุดจริง} สมมติให้เป็น $N$
    \item สร้างจำนวนใหม่ขึ้นมา: $M = N + 1$
    \item เนื่องจาก $N$ เป็นจำนวนเต็ม ดังนั้น $M$ ก็ต้องเป็นจำนวนเต็มด้วย
    \item เนื่องจาก $1 > 0$ ทำให้เราได้ว่า $M > N$
    \item จะเห็นว่า $M$ เป็นจำนวนเต็มที่มีค่ามากกว่า $N$
    \item เกิดข้อขัดแย้ง: มีจำนวนเต็ม $M$ ที่ใหญ่กว่า $N$ ซึ่ง \textcolor{Fred}{ขัดแย้งกับที่สมมติว่า $N$ ใหญ่ที่สุด}
  \end{enumerate}
  \begin{exampleblock}{สรุป}
    เกิดข้อขัดแย้ง $\Rightarrow$ ข้อสมมติที่ว่ามีจำนวนเต็มที่ใหญ่ที่สุดเป็นเท็จ $\Rightarrow$ \textbf{ไม่มีจำนวนเต็มที่ใหญ่ที่สุด}
  \end{exampleblock}
\end{frame}

\begin{frame}{ข้อขัดแย้ง (Contradiction)}
  \begin{block}{นิยาม}
    ประพจน์เชิงประกอบที่ \textbf{เป็นเท็จทุกกรณี} ไม่ว่าตัวแปรจะมีค่าใด
  \end{block}
  \vspace{0.3em}
  ตัวอย่างคลาสสิก $p \land \neg p$ (กฎแห่งการไม่ขัดแย้ง)
  \begin{center}
  \begin{tabular}{|c|c|c|}
    \hline \textbf{$p$} & \textbf{$\neg p$} & \textbf{$p\land\neg p$} \\ \hline
    \TT & \FF & \FF \\ \hline
    \FF & \TT & \FF \\ \hline
  \end{tabular}
  \end{center}
  \vspace{0.4em}
  ตัวอย่างอื่น: $\neg(p\lor\neg p)$ \quad และ \quad $(p\rightarrow q)\land(p\land\neg q)$ \\
  \textcolor{gray}{ประพจน์เดียวกันเป็นทั้งจริงและเท็จพร้อมกันไม่ได้}
\end{frame}

\begin{frame}{ประพจน์กลาง (Contingency)}
  \begin{block}{นิยาม}
    ประพจน์เชิงประกอบที่ \textbf{จริงบ้าง เท็จบ้าง} ขึ้นกับค่าของตัวแปรย่อย \\
    (ไม่เป็นทั้งสัจนิรันดร์และข้อขัดแย้ง)
  \end{block}
  \vspace{0.3em}
  ตัวอย่าง $p \land q$
  \begin{center}
  \begin{tabular}{|c|c|c|}
    \hline \textbf{$p$} & \textbf{$q$} & \textbf{$p\land q$} \\ \hline
    \TT & \TT & \TT \\ \hline
    \TT & \FF & \FF \\ \hline
    \FF & \TT & \FF \\ \hline
    \FF & \FF & \FF \\ \hline
  \end{tabular}
  \end{center}
  \vspace{0.3em}
  มีทั้ง \TT\ และ \FF\ $\Rightarrow$ เป็นประพจน์กลาง \quad (เช่น ``วันนี้ฝนตก'')
\end{frame}


%% =====================================================================
\section{ความสมมูลและกฎพีชคณิตของประพจน์}
%% =====================================================================

\begin{frame}{ความสมมูลเชิงตรรกะ (Logical Equivalence)}
  \begin{block}{นิยาม}
    ประพจน์ $P$ และ $Q$ \textbf{สมมูลกัน} เขียนแทนด้วย $P \equiv Q$ ก็ต่อเมื่อ
    มีค่าความจริงเหมือนกันทุกกรณี \\
    (หรือกล่าวคือ $P \leftrightarrow Q$เป็นสัจนิรันดร์)
  \end{block}
  \vspace{0.6em}
  วิธีตรวจสอบความสมมูลมี 2 วิธีหลัก
  \begin{enumerate}
    \item สร้าง \textbf{ตารางค่าความจริง} เปรียบเทียบทีละแถว
    \item ใช้ \textbf{กฎพีชคณิตของประพจน์} แปลงรูปจากด้านหนึ่งให้เป็นอีกด้าน
  \end{enumerate}
\end{frame}

\begin{frame}{กฎพีชคณิตของประพจน์ (1/2)}
  \begin{center}\small
  \begin{tabular}{|l|l|}
    \hline
    \textbf{ชื่อกฎ} & \textbf{รูปสมมูล} \\ \hline
    เอกลักษณ์ (Identity) & $p \land T \equiv p$, \quad $p \lor F \equiv p$ \\ \hline
    ครอบงำ (Domination) & $p \lor T \equiv T$, \quad $p \land F \equiv F$ \\ \hline
    นิจพล (Idempotent) & $p \lor p \equiv p$, \quad $p \land p \equiv p$ \\ \hline
    นิเสธซ้อน (Double Negation) & $\neg(\neg p) \equiv p$ \\ \hline
    สลับที่ (Commutative) & $p \lor q \equiv q \lor p$, \quad $p \land q \equiv q \land p$ \\ \hline
    เปลี่ยนกลุ่ม (Associative) & $(p \lor q)\lor r \equiv p \lor (q \lor r)$ \\ \hline
    แจกแจง (Distributive) & $p \lor (q \land r) \equiv (p\lor q)\land(p\lor r)$ \\ \hline
  \end{tabular}
  \end{center}
\end{frame}

\begin{frame}{กฎพีชคณิตของประพจน์ (2/2)}
  \begin{center}\small
  \begin{tabular}{|l|l|}
    \hline
    \textbf{ชื่อกฎ} & \textbf{รูปสมมูล} \\ \hline
    เดอมอร์แกน (De Morgan) & $\neg(p \land q) \equiv \neg p \lor \neg q$ \\ \hline
                            & $\neg(p \lor q) \equiv \neg p \land \neg q$ \\ \hline
    ดูดกลืน (Absorption) & $p \lor (p \land q) \equiv p$, \quad $p \land (p \lor q) \equiv p$ \\ \hline
    นิเสธ (Negation) & $p \lor \neg p \equiv T$, \quad $p \land \neg p \equiv F$ \\ \hline
    เงื่อนไข (Implication) & $p \rightarrow q \equiv \neg p \lor q$ \\ \hline
    กระจายเงื่อนไข (Exportation) & $(p\land q)\rightarrow r \equiv p \rightarrow (q \rightarrow r)$ \\ \hline
  \end{tabular}
  \end{center}
  \vspace{0.4em}
  \begin{exampleblock}{สำคัญ}
    กฎเหล่านี้คือ ``กรณีรูปธรรม'' ของกฎพีชคณิตบูลีน ที่จะใช้ออกแบบและลดรูปวงจรดิจิทัล
  \end{exampleblock}
\end{frame}

\begin{frame}{ตัวอย่าง: การลดรูปด้วยกฎพีชคณิต}
  \begin{block}{โจทย์}
    จงพิสูจน์ว่า $\neg(p \lor q) \lor (\neg p \land q) \equiv \neg p$
  \end{block}
  \vspace{0.3em}
  \textbf{วิธีทำ}
  \begin{align*}
    \neg(p \lor q) \lor (\neg p \land q)
      &\equiv (\neg p \land \neg q) \lor (\neg p \land q) && \text{(เดอมอร์แกน)} \\
      &\equiv \neg p \land (\neg q \lor q) && \text{(แจกแจง)} \\
      &\equiv \neg p \land T && \text{(นิเสธ)} \\
      &\equiv \neg p && \text{(เอกลักษณ์)}
  \end{align*}
  \begin{exampleblock}{คำตอบ}
    $\neg(p \lor q) \lor (\neg p \land q) \equiv \neg p$
  \end{exampleblock}
\end{frame}




%% =====================================================================
\section{กิจกรรมในชั้นเรียน (แบบฝึกหัดทำในห้อง)}
%% =====================================================================

\begin{frame}{วิธีทำกิจกรรม}
  \begin{block}{คำชี้แจง}
    ให้นักศึกษาทำกิจกรรมต่อไปนี้ \textbf{ในชั้นเรียน} ลงในสมุด/กระดาษ
    ภายในเวลาที่กำหนด แล้วร่วมเฉลยและอภิปรายพร้อมกัน
  \end{block}
  \vspace{0.5em}
  \begin{itemize}
    \item ทำเดี่ยวหรือจับคู่ก็ได้ (ตามที่อาจารย์กำหนด)
    \item เมื่อหมดเวลา อาจารย์จะคลิกเพื่อ \textbf{เฉลย} ทีละข้อ
    \item อาสาสมัครออกมาอธิบายวิธีคิดหน้าชั้น
  \end{itemize}
  \vspace{0.5em}
  \begin{exampleblock}{เกณฑ์}
    เน้นกระบวนการคิดและการให้เหตุผล ไม่ใช่แค่คำตอบสุดท้าย
  \end{exampleblock}
\end{frame}

\begin{frame}{กิจกรรมที่ 1: จำแนกประเภทประพจน์ \eng{6 นาที}}
  จงจำแนกว่าแต่ละประพจน์เป็น \textbf{สัจนิรันดร์ / ข้อขัดแย้ง / ประพจน์กลาง}
  \begin{enumerate}
    \item $p \lor \neg p$
    \item $p \land \neg p$
    \item $(p \rightarrow q) \rightarrow (\neg q \rightarrow \neg p)$
    \item $p \rightarrow q$
  \end{enumerate}
  \pause
  \begin{exampleblock}{เฉลย}
    \footnotesize
    1) \textcolor{Tgreen}{สัจนิรันดร์} \quad 2) \textcolor{Fred}{ข้อขัดแย้ง} \quad
    3) \textcolor{Tgreen}{สัจนิรันดร์} (เพราะ $p\rightarrow q \equiv \neg q\rightarrow\neg p$ จึงเป็น $X\rightarrow X$) \quad
    4) \textcolor{navy}{ประพจน์กลาง}
  \end{exampleblock}
\end{frame}

\begin{frame}{กิจกรรมที่ 2: กฎเดอมอร์แกน \eng{6 นาที}}
  จงเขียนนิเสธของประพจน์ต่อไปนี้ให้อยู่ในรูปที่ ``นิเสธอยู่ติดตัวแปร'' โดยใช้กฎเดอมอร์แกน
  \begin{enumerate}
    \item $\neg(p \land \neg q)$
    \item $\neg(\neg p \lor q)$
    \item $\neg(p \lor (q \land r))$
  \end{enumerate}
  \pause
  \begin{exampleblock}{เฉลย}
    \footnotesize
    1) $\neg p \lor q$ \qquad
    2) $p \land \neg q$ \qquad
    3) $\neg p \land (\neg q \lor \neg r)$
  \end{exampleblock}
\end{frame}

\begin{frame}{กิจกรรมที่ 3: ตรวจสอบความสมมูล \eng{8 นาที}}
  จงสร้างตารางค่าความจริงเพื่อตรวจว่า \quad
  $p \rightarrow (q \land r)$ \quad สมมูลกับ \quad $(p\rightarrow q)\land(p\rightarrow r)$ \quad หรือไม่
  \pause
  \begin{center}\scriptsize
  \begin{tabular}{|c|c|c|c|c|c|c|c|}
    \hline
    $p$ & $q$ & $r$ & $q\land r$ & $p\rightarrow(q\land r)$ & $p\rightarrow q$ & $p\rightarrow r$ & $(p\rightarrow q)\land(p\rightarrow r)$ \\ \hline
    \TT & \TT & \TT & \TT & \TT & \TT & \TT & \TT \\ \hline
    \TT & \TT & \FF & \FF & \FF & \TT & \FF & \FF \\ \hline
    \TT & \FF & \TT & \FF & \FF & \FF & \TT & \FF \\ \hline
    \TT & \FF & \FF & \FF & \FF & \FF & \FF & \FF \\ \hline
    \FF & \TT & \TT & \TT & \TT & \TT & \TT & \TT \\ \hline
    \FF & \TT & \FF & \FF & \TT & \TT & \TT & \TT \\ \hline
    \FF & \FF & \TT & \FF & \TT & \TT & \TT & \TT \\ \hline
    \FF & \FF & \FF & \FF & \TT & \TT & \TT & \TT \\ \hline
  \end{tabular}
  \end{center}
  \footnotesize สองคอลัมน์สุดท้ายเท่ากันทุกแถว $\Rightarrow$ \textbf{\textcolor{Tgreen}{สมมูลกัน}}
\end{frame}


%% ====================== สรุป + แบบฝึกหัด ======================
\begin{frame}{สรุปสัปดาห์ที่ 2}
  \begin{itemize}
    \item \textbf{สัจนิรันดร์ (Tautology)}: จริงทุกกรณี (เช่น $p \lor \neg p$, $(p \land q) \rightarrow p$)
    \item \textbf{ข้อขัดแย้ง (Contradiction)}: เท็จทุกกรณี (เช่น $p \land \neg p$)
    \item \textbf{ประพจน์กลาง (Contingency)}: จริงบ้างเท็จบ้าง (เช่น $p \rightarrow q$)
    \item \textbf{ความสมมูล ($P \equiv Q$)}: ค่าความจริงเหมือนกันกรณีต่อกรณี
    \item \textbf{กฎพีชคณิตของประพจน์}: กฎสลับที่ เปลี่ยนกลุ่ม แจกแจง ดูดกลืน และเดอมอร์แกน
  \end{itemize}
  \begin{exampleblock}{}
    \centering สัปดาห์ถัดไป: ตัวเชื่อมพิเศษ (XOR, NAND, NOR) ตัวบ่งปริมาณ และกฎการอนุมาน
  \end{exampleblock}
\end{frame}

\begin{frame}{แบบฝึกหัดมอบหมาย (สัปดาห์ที่ 2)}
  \begin{enumerate}\small
    \item จำแนกประพจน์เป็นสัจนิรันดร์/ข้อขัดแย้ง/ประพจน์กลาง:
          \begin{itemize}
            \item $(p\rightarrow q)\lor(q\rightarrow p)$
            \item $\neg(p\lor\neg p)$
            \item $(p\land q)\rightarrow p$
          \end{itemize}
    \item ตรวจสอบความสมมูลเชิงตรรกะด้วยกฎพีชคณิตประพจน์:
          \begin{itemize}
            \item $\neg(p\land q)$ กับ $\neg p\lor\neg q$
          \end{itemize}
  \end{enumerate}
\end{frame}

\begin{frame}[plain]
  \vfill
  \centering
  {\Large\color{navy}\textbf{จบสัปดาห์ที่ 2}}\par
  \vspace{0.8em}
  {\color{gray} สัปดาห์ถัดไป: สัปดาห์ที่ 3 · ตัวเชื่อมเสริม ตัวบ่งปริมาณ และการอนุมาน}
  \vfill
\end{frame}
